\section{Analysis}
\label{sec:analysis}

\subsection{What the Controlled Results Support}

The controlled tables support a narrow mechanism claim. Directly using
wavelet-derived local response as an anomaly map is not sufficient, so the
method should not be described as local-cue fusion. Semantic-only prototype
adaptation is already strong, which means that much of the gain comes from
image-conditioned prototype estimation itself. The remaining improvement from
boundary-aware reliability is smaller but consistent across all four reported
metrics on both datasets. This is the level of claim supported by the current
evidence.

\subsection{Why Pixel AUPRO Is Informative}

The largest controlled margin is usually in Pixel AUPRO. This is consistent
with the proposed mechanism: a better image-conditioned reference should reduce
region-level false positives and improve localization quality, even when the
absolute Pixel AUROC change is modest. The normal-image stability table gives a
second check in the same direction: conservative calibration does not increase
the area marked abnormal on normal images.

\subsection{What the Results Do Not Support}

The current data do not support the stronger claim that frequency alone solves
zero-shot anomaly localization. They also do not support presenting external
baseline tables as strict same-protocol rankings unless the split, backbone,
input resolution, post-processing, training data, and metric implementation are
verified for each row. External baselines can define the research context, but
the mechanism claim in this draft is grounded in the controlled comparisons:
direct fusion, semantic-only adaptation, reliability design, conservative
calibration, and normal-image stability.
